\chapter{SSR Authoring (v1): Rendering Blocks to HTML Strings}

Server-side rendering (SSR) is how you produce initial HTML on the server so the browser can show an app quickly.

In HRBR v1, SSR is intentionally simple and primarily exists to support tests and examples.
The important property is not performance; it’s \textbf{compatibility} with hydration.

This chapter is grounded in:

\begin{itemize}
  \item \texttt{runtime/ssr.ts}
  \item \texttt{runtime/__tests__/ssr.test.ts}
  \item and the follow-on behavior in \texttt{runtime/hydration.ts} (Chapter 15).
\end{itemize}

\section{A tiny contract (what SSR v1 guarantees)}

SSR is not “render an app” in the abstract. In HRBR v1 it is narrowly scoped:

\begin{enumerate}
  \item You provide a \texttt{BlockDef} (template + slots).
  \item You optionally provide \texttt{initialValues} (a slot-key \(\rightarrow\) value map).
  \item The runtime returns an HTML string that matches the same template structure the client will mount/hydrate.
\end{enumerate}

Put differently:

\begin{quote}
SSR v1 is \emph{mountBlock + serialize}. Hydration v1 is \emph{validate + attach, else remount}.
\end{quote}

This tight coupling is deliberate. It makes SSR and hydration “round-trip compatible” by construction.

\section{The contract: \texttt{renderBlockToString(def, options?) -> string}}

The only exported function is:

\begin{verbatim}
renderBlockToString(def: BlockDef, options?: { initialValues?: Record<string, unknown> }): string
\end{verbatim}

The options only support one thing today: \texttt{initialValues}.

This mirrors the core block update contract:

\begin{quote}
A block update is just a map \texttt{\{ [slotKey]: value \}}.
\end{quote}

SSR is simply the "apply values, then serialize" version of that idea.

One way to think about it:

\begin{itemize}
  \item \textbf{Client mount:} template \(\rightarrow\) DOM instance \(+\) slot node references.
  \item \textbf{Client update:} slot values \(\rightarrow\) DOM writes.
  \item \textbf{Server render:} slot values \(\rightarrow\) DOM writes \(\rightarrow\) serialize.
\end{itemize}

\section{Implementation: DOM-based SSR for correctness}

\texttt{runtime/ssr.ts} implements SSR using the platform DOM:

\begin{enumerate}
  \item create a detached host element: \texttt{document.createElement('div')}
  \item mount the block into that host via \texttt{mountBlock(def, host, initialValues)}
  \item read \texttt{host.innerHTML}
  \item destroy the mounted block to avoid leaks
  \item return the HTML string
\end{enumerate}

This is documented directly in the function comment:

\begin{itemize}
  \item it uses a detached DOM
  \item it serializes via \texttt{innerHTML}
  \item it’s primarily intended for tests/examples
  \item a future version could emit HTML without a DOM
\end{itemize}

The crucial correctness choice is that SSR applies \texttt{initialValues} by calling \texttt{mountBlock(def, host, initialValues)}.
That means:

\begin{itemize}
  \item all the slot semantics from Chapter 14 apply (boolean attributes, SVG namespaces, input properties, etc.)
  \item the generated markup shape matches what the client expects to attach to
\end{itemize}

\subsection{Why DOM-based SSR is acceptable in v1}

For a production SSR renderer, involving a full DOM implementation server-side can be expensive.

But for v1 HRBR, this has two big advantages:

\begin{itemize}
  \item it is hard to get HTML escaping and attribute semantics wrong
  \item it produces markup exactly consistent with the same block runtime used in the client
\end{itemize}

So correctness comes for free.

There’s also a subtle lifecycle advantage: because SSR uses \texttt{mountBlock}, tearing down is just \texttt{mounted.destroy()}.
That ensures listeners are removed and the temporary root is cleaned up, even in a test environment.

\section{SSR must round-trip with hydration}

SSR isn't useful unless the client can hydrate it.

Hydration v1 makes strong assumptions (Chapter 15):

\begin{itemize}
  \item the root element tagName matches \texttt{templateHTML}
  \item slot paths resolve to the expected node types
  \item element slots match expected tagNames along the path
\end{itemize}

Therefore, SSR output must preserve the template’s structure.
Because SSR uses \texttt{mountBlock()}, it naturally meets the same invariant:

\begin{quote}
SSR uses the same template parser and slot patching semantics as client mount.
\end{quote}

This is the important “authoring rule” for SSR v1:

\begin{quote}
If you change \texttt{templateHTML} shape, you must change both server render and client hydrate together.
\end{quote}

You should treat \texttt{templateHTML} as a stable ABI between server and client.

\section{The end-to-end spec: render \(\rightarrow\) hydrate \(\rightarrow\) update}

The key test is \texttt{runtime/__tests__/ssr.test.ts}:

\begin{enumerate}
  \item Define a block with \texttt{templateHTML} and a text slot.
  \item Render HTML with \texttt{renderBlockToString(block, \{ initialValues \})}.
  \item Insert the returned string into the DOM.
  \item Call \texttt{hydrateBlock(block, host)}.
  \item Verify DOM identity is preserved (same root, same text node).
  \item Mount a compiled block via \texttt{mountCompiledBlock(...)} and update.
\end{enumerate}

This accomplishes two things:

\begin{itemize}
  \item it proves that SSR output is "hydrate compatible"
  \item it shows how hydration and compiled blocks can be composed (a realistic client story)
\end{itemize}

One implementation-level nuance in the test is worth calling out.
After updating the signal with \texttt{setCount('B')}, it also calls:

\begin{quote}
	exttt{mounted.update(\{ value: 'B' \})}
\end{quote}

That makes the assertion deterministic and keeps the test independent of scheduler timing.
It’s modeling the fact that \texttt{mountCompiledBlock} can schedule work, and test environments often differ in how timers/microtasks flush.

\section{Limitations and the intended direction}

The HRBR whitepaper explicitly lists a non-goal for v1:

\begin{quote}
A DOM-free string renderer (v1 SSR uses a DOM-based approach for tests/examples)
\end{quote}

A future SSR implementation could:

\begin{itemize}
  \item parse \texttt{templateHTML} once into a token stream
  \item apply slot writes into a string builder
  \item avoid constructing DOM nodes entirely
\end{itemize}

Even if SSR becomes DOM-free later, the invariants from hydration still apply:

\begin{itemize}
  \item the root tag must match
  \item the structure along slot paths must match
  \item text nodes must exist where text slots expect them
\end{itemize}

But even then, the contract remains:

\begin{quote}
SSR output must match the template structure so hydration can do a no-rewrite attach.
\end{quote}

\section{Takeaways}

\begin{itemize}
  \item SSR v1 is intentionally minimal: mount a block in a detached DOM and serialize \texttt{innerHTML}.
  \item \texttt{initialValues} uses the same slot update contract as client updates.
  \item The primary goal is hydration compatibility, not server performance.
  \item The end-to-end test shows SSR + hydration + compiled updates working together.
\end{itemize}

\section{Walking the tests (the executable spec)}

SSR v1 has one test file, but it covers the entire story.

\subsection{\texttt{runtime/__tests__/ssr.test.ts}}

The test \texttt{renders a block to HTML string, hydrates it, and updates via mountCompiledBlock} enforces these guarantees:

\begin{itemize}
  \item \textbf{SSR applies initial values:} \texttt{renderBlockToString(..., \{ initialValues \})} must produce HTML containing \texttt{'A'}.
  \item \textbf{Hydration preserves identity:} after \texttt{hydrateBlock}, the root element and the text node under \texttt{#t} are the same objects as before hydration.
  \item \textbf{Compiled updates work on hydrated DOM:} wiring a signal through \texttt{mountCompiledBlock} and then updating should change text to \texttt{'B'}.
  \item \textbf{Cleanup is explicit:} the test disposes the compiled mount and destroys the hydrated block.
\end{itemize}

Because this file touches SSR, hydration, signals, and compiled blocks, it acts as a integration-level "safety rail" for the whole pipeline.

\section{Online resources}

\begin{itemize}
  \item HTML serialization via DOM:
  \href{https://developer.mozilla.org/en-US/docs/Web/API/Element/innerHTML}{MDN: \texttt{Element.innerHTML}}.
  \item Server rendering and hydration overview:
  \href{https://web.dev/rendering-on-the-web/}{web.dev: Rendering on the Web}.
  \item Template element (used by mounting/hydration):
  \href{https://developer.mozilla.org/en-US/docs/Web/HTML/Element/template}{MDN: \texttt{<template>}}.
  \item A compiler-first argument for skipping VDOM in the steady state:
  \href{https://svelte.dev/blog/virtual-dom-is-pure-overhead}{Svelte: “Virtual DOM is pure overhead”}.
\end{itemize}
